\begin{table}[H]
\centering
\caption{Multi-Cutoff RDD: Effect of Monitoring Intensity on Violations}
\begin{threeparttable}
\begin{tabular}{lcccc}
\toprule
Outcome & Estimate & Bandwidth & $N$ (left/right) & $p$-value \\
\midrule
Any coliform MCL violation & -0.0109 & 329 & 6,319 / 4,848 & 0.373 \\
 & (0.0128) & & & \\
Any health-based violation & -0.0019 & 463 & 8,915 / 6,180 & 0.944 \\
 & (0.0161) & & & \\
Any non-coliform MCL violation & -0.0204* & 368 & 7,146 / 5,278 & 0.057 \\
 & (0.0119) & & & \\
Number of coliform violations & -0.1976 & 392 & 7,590 / 5,528 & 0.175 \\
 & (0.1677) & & & \\
\bottomrule
\end{tabular}
\begin{tablenotes}[flushleft]
\small
\item \textit{Notes:} Each row reports a separate local linear RDD. The running variable is population served, normalized as distance to the nearest of nine monitoring thresholds (1,000; 2,500; 3,300; 4,100; 4,900; 5,800; 6,700; 7,600; 8,500). Triangular kernel with MSE-optimal bandwidth \citep{calonico2014robust}. Bias-corrected robust standard errors in parentheses and $p$-values. * $p<0.10$, ** $p<0.05$, *** $p<0.01$.
\end{tablenotes}
\end{threeparttable}
\label{tab:main}
\end{table}
